\documentclass[11pt]{article}
\usepackage[margin=1in]{geometry}
\usepackage{amsmath,amssymb,amsthm}
\usepackage{graphicx}
\usepackage{hyperref}
\usepackage{enumitem}

\newtheorem{theorem}{Theorem}
\newtheorem{lemma}{Lemma}
\newtheorem{remark}{Remark}

\newcommand{\R}{\mathbb{R}}
\newcommand{\N}{\mathbb{N}}
\newcommand{\Ball}{B}
\newcommand{\ext}{\operatorname{ext}}
\newcommand{\cconv}{\overline{\operatorname{co}}}
\newcommand{\DeltaX}{\Delta_X}

\title{A Counterexample to the BD-Polyhedral Characterisation of convex-DLD2P}
\author{Run \texttt{fa\_banach\_001}, lane 9}
\date{June 18, 2026}

\begin{document}
\maketitle

\section*{Status}

Candidate counterexample, likely valid.  The packet answers negatively
Problem 3.12 of Martin and Rueda Zoca~\cite{MRZ2010}, which asks whether a
real \(L_1\)-predual has the convex diametral local diameter two property
(convex-DLD2P) if and only if it fails to be (BD) polyhedral.

\section*{Source Problem}

The problem appears on page 12 of the source paper.

\begin{center}
\includegraphics[width=\linewidth]{figures/open_problem_crop.png}
\end{center}

In the notation of the source, \(\Delta_X\) denotes the set of \(\Delta\)-points
of \(S_X\), and \(X\) has the convex-DLD2P when
\[
  \Ball_X=\cconv(\Delta_X).
\]
For a real Banach space \(X\), property (BD) is
\[
  \sup\{x^*(x): x^*\in \ext(\Ball_{X^*})\setminus D(x)\}<1
  \quad (x\in S_X),
\]
where \(D(x)=\{x^*\in S_{X^*}: x^*(x)=1\}\).

\section*{Counterexample}

Let \(c\) be the real Banach space of convergent scalar sequences with the
supremum norm.  For \(x=(x_n)_{n\geq 1}\in c\), write
\[
  x_0=\lim_{n\to\infty}x_n.
\]
Define
\[
  X=\{x\in c: 2x_0=x_1+x_2\}.
\]

\section*{Proof Intuition}

The source theorem characterises \(\Delta\)-points in an \(L_1\)-predual by
looking at how infinitely many extreme dual functionals evaluate on a point.
For this hyperplane the extreme dual functionals are just the coordinate
functionals, so a point is a \(\Delta\)-point exactly when its limiting
coordinate has modulus one.  But the equation \(2x_0=x_1+x_2\) then forces
\(x_1=x_2=x_0\).  Hence all \(\Delta\)-points live in the closed convex slice
\(x_1=x_2=x_0\), while the unit ball also contains points with \(x_1\neq x_2\).
At the same time, the point whose first two coordinates and limit are \(1\),
but whose tail approaches \(1\) from below, witnesses failure of (BD).

\begin{theorem}
The real Banach space \(X=\{x\in c:2\lim_n x_n=x_1+x_2\}\) is an
\(L_1\)-predual, fails to be (BD) polyhedral, and does not have the
convex-DLD2P.
\end{theorem}

\begin{proof}
First we record that \(X\) is an \(L_1\)-predual.  In the notation of
Casini, Miglierina, and Piasecki~\cite{CMP2014}, let
\[
  f=\left(\frac12,-\frac14,-\frac14,0,0,\ldots\right)\in \ell_1,
\]
where the first coordinate acts on \(x_0=\lim_n x_n\), and the following
coordinates act on \(x_1,x_2,\ldots\).  Then \(\|f\|_1=1\) and
\[
  \ker f=\{x\in c:\tfrac12x_0-\tfrac14x_1-\tfrac14x_2=0\}=X.
\]
Since \(|f_1|=1/2\) and \(|f_j|<1/2\) for every \(j\geq 2\), the
hyperplane theorem of~\cite[Theorem 4.3]{CMP2014} applies: the map
\[
  \phi:\ell_1\longrightarrow X^*,\qquad
  \phi(a)(x)=\sum_{n=1}^\infty a_nx_n,
\]
is an onto isometry.  Consequently \(X^*\) is isometric to \(\ell_1\), and
\[
  \ext(\Ball_{X^*})=\{\pm e_n^*:n\in\N\},
\]
where \(e_n^*(x)=x_n\).

We next identify the \(\Delta\)-points of \(X\).  By
\cite[Theorem 3.2]{MRZ2010}, for a real \(L_1\)-predual \(Y\) and
\(u\in S_Y\), \(u\) is a \(\Delta\)-point if and only if, for every
\(\varepsilon>0\), the set
\[
  \{e^*\in \ext(\Ball_{Y^*}): |e^*(u)|>1-\varepsilon\}
\]
contains infinitely many pairwise linearly independent elements.  In our
space this says that \(u\in S_X\) is a \(\Delta\)-point if and only if
infinitely many coordinates \(u_n\) are arbitrarily close in modulus to
\(1\).  Because \(u\in c\), this is equivalent to
\[
  |u_0|=1.
\]
If \(|u_0|=1\), the constraint \(2u_0=u_1+u_2\) and
\(|u_1|,|u_2|\leq 1\) force \(u_1=u_2=u_0\).  Hence
\[
  \DeltaX\subset M:=\{x\in \Ball_X:x_1=x_2=x_0\}.
\]
The set \(M\) is closed and convex, so \(\cconv(\DeltaX)\subset M\).
However,
\[
  z=(1,-1,0,0,\ldots)\in c,\qquad z_0=0,
\]
belongs to \(X\), since \(2z_0=0=z_1+z_2\), and \(\|z\|=1\).  But
\(z\notin M\).  Therefore
\[
  \Ball_X\neq \cconv(\DeltaX),
\]
so \(X\) does not have the convex-DLD2P.

It remains to check failure of (BD).  Let \(u\in X\) be given by
\[
  u_0=1,\qquad u_1=u_2=1,\qquad u_n=1-\frac1n\quad(n\geq 3).
\]
Then \(u\in S_X\).  For every \(n\geq 3\),
\[
  e_n^*(u)=1-\frac1n<1,
\]
so \(e_n^*\notin D(u)\).  Nevertheless \(e_n^*(u)\to 1\).  Thus
\[
  \sup\{x^*(u):x^*\in\ext(\Ball_{X^*})\setminus D(u)\}=1,
\]
which is exactly the failure of (BD) at \(u\).
\end{proof}

\section*{Verification Notes}

The proof uses only two external facts:
\begin{enumerate}[label=(\arabic*)]
\item the hyperplane duality theorem of Casini--Miglierina--Piasecki for
      \(W_f\subset c\), applied to \(f=(1/2,-1/4,-1/4,0,\ldots)\);
\item Theorem 3.2 of the source paper, characterising \(\Delta\)-points in
      \(L_1\)-preduals by extreme dual functionals.
\end{enumerate}
The script \texttt{code/verify\_counterexample.py} checks the displayed
hyperplane equations and the elementary witnesses, but it is not used as a
substitute for the proof.

\section*{Novelty Check}

Before promotion, the run indexes were searched for
\(\texttt{2010.15936}\), \(\texttt{convex-DLD2P}\),
\(\texttt{BD polyhedral}\), \(\texttt{Problem 3.12}\), and the explicit
hyperplane pattern.  No duplicate packet or attempt was found.

Web searches on June 18, 2026 for exact and close phrases, including
``convex-DLD2P'' with ``BD polyhedral'', ``Problem 3.12'' with
``convex-DLD2P'', and the explicit hyperplane pattern \(2\lim x=x_1+x_2\),
did not surface a later paper answering the problem or this construction.

\section*{Human Review Recommendation}

The main review point is to confirm the exact indexing convention in
\cite[Theorem 4.3]{CMP2014}.  With the convention used there, the first
coordinate of \(f\) multiplies \(x_0=\lim_n x_n\), and the example has
\(|f_1|=1/2\), so the cited theorem applies.

\begin{thebibliography}{9}

\bibitem{MRZ2010}
M. Martin and A. Rueda Zoca,
\emph{Daugavet property in projective symmetric tensor products of Banach spaces},
arXiv:2010.15936.

\bibitem{CMP2014}
E. Casini, E. Miglierina, and L. Piasecki,
\emph{Hyperplanes in the space of convergent sequences and preduals of \(\ell_1\)},
arXiv:1410.7801.

\end{thebibliography}

\end{document}
